\begin{figure}[H]
    \centering
\begin{tikzpicture}[
    scale=2,
    cota/.style={<->, >=stealth, thin, gray!80!black},
    retangulo/.style={draw=gray!80!black, thick},
    divisao/.style={draw=black, thick, dashed},          % Charneiras / linhas de divisão
    eixo/.style={->, >=stealth, thick, gray!80!black},
    corte/.style={draw=black, dash dot, thick},          % Linhas dos cortes A e B
    setaCorte/.style={->, >=stealth, thick, black}
]

    % 1. Coordenadas dos Vértices do Retângulo
    \coordinate (A) at (0, 0);     % Canto inferior esquerdo
    \coordinate (B) at (3, 0);     % Canto inferior direito
    \coordinate (C) at (3, 2);     % Canto superior direito
    \coordinate (D) at (0, 2);     % Canto superior esquerdo

    % Interseções internas das linhas de divisão (charneiras plásticas)
    \coordinate (P1) at (0.732, 1.268);
    \coordinate (P2) at (1.732, 1.268);

    % 2. Hachuras externas sem borda (Engastes: base e lateral direita)
    \path[pattern={Lines[angle=45, distance=6pt, line width=0.5pt]}, pattern color=gray!60] (0, -0.3) rectangle (3.3, 0);
    \path[pattern={Lines[angle=45, distance=6pt, line width=0.5pt]}, pattern color=gray!60] (3, 0) rectangle (3.3, 2);

    % 3. Desenho do Retângulo Principal
    \draw[retangulo] (A) rectangle (C);

    % 4. Linhas de Divisão de Carga (Charneiras)
    \draw[divisao] (A) -- (P1);
    \draw[divisao] (D) -- (P1);
    \draw[divisao] (P1) -- (P2);
    \draw[divisao] (B) -- (P2);
    \draw[divisao] (C) -- (P2);

    % 5. Linhas dos Cortes A-A e B-B
    % Corte A-A (Horizontal pelo centro em y = 1.0)
    \draw[corte] (-0.5, 1.0) -- (3.5, 1.0);
    \draw[setaCorte] (-0.5, 1.0) -- (-0.5, 1.3) node[above, font=\small\bfseries] {A};
    \draw[setaCorte] (3.5, 1.0) -- (3.5, 1.3) node[above, font=\small\bfseries] {A};

    % Corte B-B (Vertical pelo centro em x = 1.5)
    \draw[corte] (1.5, -0.7) -- (1.5, 2.5);
    \draw[setaCorte] (1.5, 2.5) -- (1.8, 2.5) node[right, font=\small\bfseries] {B};
    \draw[setaCorte] (1.5, -0.7) -- (1.8, -0.7) node[right, font=\small\bfseries] {B};

    % 6. Textos das divisões (1, 2, 3, 4)
    \node[font=\small\bfseries, gray!90!black, fill=white, inner sep=1.5pt, circle] at (0.35, 1.0) {$A_y$};
    \node[font=\small\bfseries, gray!90!black, fill=white, inner sep=1.5pt, circle] at (1.23, 1.6) {$A_x$};
    \node[font=\small\bfseries, gray!90!black, fill=white, inner sep=1.5pt, circle] at (2.4, 1.0) {$A'_y$};
    \node[font=\small\bfseries, gray!90!black, fill=white, inner sep=1.5pt, circle] at (1.23, 0.5) {$A'_x$};

    % 7. Arcos e Rótulos dos Ângulos
    % Canto A (inferior esquerdo) - 60°
    \draw[thin, gray!80!black] (0.6, 0) arc (0:60:0.6);
    \node[font=\tiny, gray!90!black] at (0.75, 0.3) {$60^\circ$};

    % Canto D (superior esquerdo) - 45° EM RELAÇÃO À BORDA ESQUERDA (VERTICAL)
    \draw[thin, gray!80!black] (0, 1.5) arc (270:315:0.5);
    \node[font=\tiny, gray!90!black] at (0.3, 1.45) {$45^\circ$};

    % Canto B (inferior direito) - 45°
    \draw[thin, gray!80!black] (2.5, 0) arc (180:135:0.5);
    \node[font=\tiny, gray!90!black] at (2.35, 0.25) {$45^\circ$};

    % Canto C (superior direito) - 60°
    \draw[thin, gray!80!black] (3, 1.5) arc (270:210:0.5);
    \node[font=\tiny, gray!90!black] at (2.7, 1.4) {$60^\circ$};

    % 8. Eixos de Coordenadas x e y
    \draw[eixo] (0.1, 0.1) -- (0.5, 0.1) node[above=-1pt, font=\tiny\bfseries] {$x$};
    \draw[eixo] (0.1, 0.1) -- (0.1, 0.5) node[right=-1pt, font=\tiny\bfseries] {$y$};

    % 9. Cotas de Dimensão
    %\def\offsetH{1.1}
    %%\draw[cota] (0, -\offsetH) -- (3, -\offsetH) 
     %   node[midway, below=3pt, font=\small] {$476\text{ cm}$};

    %\def\offsetV{0.8}
    %%\draw[cota] (-\offsetV, 0) -- (-\offsetV, 2) 
     %   node[midway, left=3pt, font=\small, rotate=90, anchor=south] {$256 \text{ cm}$};

\end{tikzpicture}
    %\caption{Divisão de áreas de carga de laje com linhas de corte A-A e B-B}
    %\label{fig:divisao_laje}
\end{figure}